\documentclass[11pt,a4paper,twocolumn]{article}

% --- Grundlegende Pakete ---
\usepackage[utf8]{inputenc}
\usepackage[T1]{fontenc}
\usepackage[german]{babel}
\usepackage{amsmath, amssymb, amsthm} % Mathematische Symbole
\usepackage{bm}                       % Fette Mathe-Symbole
\usepackage{graphicx}
\usepackage{xcolor}

% --- PGFPlots für Ihre SageMath-Daten ---
\usepackage{pgfplots}
\pgfplotsset{compat=1.18}
\usepgfplotslibrary{statistics}

% --- Definitionen für die #Energiedoku ---
\newcommand{\Tc}{T_{\text{c}}}
\newcommand{\OmegaH}{\mathcal{P}_{\mathbb{H},\mathrm{fluct}}}

% --- Metadaten ---
\title{Arithmetische Holonomie und spektrale Universalität am kritischen Punkt $\Tc$}
\author{Thomas Hofbauer \\ \small Bamberg, Deutschland}
\date{Mai 2026}

\begin{document}

\maketitle

\begin{abstract}
Dieses Manuskript dokumentiert die statistische Kongruenz zwischen der quaternionischen Bamberger Zeitwürfel-Sonde und der Riemannschen Zeta-Funktion. Mittels Kolmogorov-Smirnov-Validierung wird nachgewiesen, dass am Punkt $\Tc \approx \pi$ eine Universalitätsklasse erreicht wird.
\end{abstract}

\section{Der Fluktuationsoperator}
Der entscheidende Fortschritt für den Vergleich ist der Fluktuationsoperator:
\begin{equation}
    \OmegaH(\Omega;T) = \left| \sum_k q_k(\tau_k-1)e^{-i\Omega k} \right|^2
\end{equation}
Hierbei isolieren wir den Jitter $(\tau_k-1)$, um die spektrale Steifigkeit des Systems zu messen.

\section{Statistische Validierung}
Die numerische Analyse liefert am kritischen Punkt eine KS-Statistik von $D \approx 0,048$ bei einem $p$-Wert von $\approx 0,199$. Dies bestätigt die Relation:
\begin{equation}
    \OmegaH(\Omega;\Tc) \sim P_\zeta(\Omega)
\end{equation}

% Hier binden Sie Ihre exportierten .dat Dateien ein:
% \begin{figure}
% \begin{tikzpicture}
% \begin{axis}[xlabel=$\Omega$, ylabel=$P(\Omega)$]
%     \addplot table {sonde_tc_spectrum.dat};
% \end{axis}
% \end{tikzpicture}
% \end{figure}

\end{document}